\chapter{The Base Model}
\label{ch:m0}
The first model is very simple and the unusual elements (like the unemployment benefits) does not add complexity to it.

\begin{table}
    \centering
    \begin{tabular}{l|ccc|l}
                 & $\mathcal{H}$             & $\mathcal{F}$               & $\mathcal{G}$    & $\Sigma$ \\
        \midrule
        Money    & $+\mathbf{M}_\mathcal{H}$ & $+\mathbf{M}_{\mathcal{F}}$ & $-\mathbf{M}$    & 0        \\
        \midrule
        $\Sigma$ & $+V_\mathcal{H}$          & $+V_{\mathcal{F}}$          & $+V_\mathcal{G}$ & 0        \\
    \end{tabular}
    \caption{Balance sheet for the Base model}
    \label{tab:bs}
\end{table}

As shown in the balance sheet (matrix \ref{tab:bs}) for the model there is only one financial asset (money) issued by the consolidated public sector (including both the government and the central bank) as a debt for itself and a credit for the other sector, which bearing no interests.

Additionally, there is one real homogenous good consumed by both the households and the government that is depleted at the end of each period.
As such there is no accumulation of the good between periods, and it is not relevant to include it in the balance sheet (which represent the state of the system at the end of each period).

\begin{table}
    \centering
    \begin{tabular}{l|ccc|l}

                       & $\mathcal{H}$                   & $\mathcal{F} $                  & $\mathcal{G}$       & $\Sigma$ \\
        \midrule
        Consumption    & $-p C$                          & $+p C$                          &                     & 0        \\
        Gvt Spending   &                                 & $+p G$                          & $-p G$              & 0        \\
        Benefits       & $+UB$                           &                                 & $-UB$               & 0        \\
        Wages          & $+W$                            & $-W$                            &                     & 0        \\
        Profits        & $+P$                            & $-P$                            &                     & 0        \\
        Taxes          &
        $-T$           &                                 & $+T$                            & 0                              \\
        \midrule
        $\Delta$ Money & $-\Delta\mathbf{M}_\mathcal{H}$ & $-\Delta\mathbf{M}_\mathcal{F}$ & $+\Delta\mathbf{M}$ & 0        \\
        \midrule
        $\Sigma$       & 0                               & 0                               & 0                   & 0        \\
    \end{tabular}
    \caption{Transactions and Flows matrix for the Base model}
    \label{tab:fof}
\end{table}

The flow of funds matrix (\ref{tab:fof}) summarizes the events of the model, with the associated flows of value.
First two rows (Consumption and Government Spending) relates to the goods market: the homogenous good is sold to both the households and the government.
The third row (Benefits) indicates the unemployment benefits paid by the government to the households that do not receive a wage.
The fourth row (Wages) summarizes the labour market, in which firms employ households and pay them a wage.
The fifth and sixth rows (Profits and Taxes) clean up the time-step: firms transfer to the households their remaining liquidity as distributed profits, and households pay taxes on wages, consumption, and profits.

\section{The aggregate version}
\subsection{Behavioural equations}
\label{sec:behave}
The behavioural equations of the aggregate version are written such there is no simultaneity among equations, making it easier to extend to the agent based version later: while the set of equation of an SFC can be solved iteratively for each time-step, ignoring the order and allowing for simultaneity and interdependency among equations, and refining the result on each iteration until the system converges, the same strategy is not feasible for an ABM.

The feasibility of the iterative method is due the low number of variables that are preserved from the previous time period as input for the following one and the relative low number of operations implied by the equations: in an ABM each variable is replicated for each agent, each equation is repeated for each agent and, often, the necessity to disaggregate the behaviour requires introducing many operations, often including choices, random events and conditional checks.
This increment the time to solve a single iteration of the ABM time-step of some magnitudes, making unfeasible to iterate multiple times to solve a single time-step.
As a consequence the sequence of the events needs to be clearly defined in the AB version to be sure that each equation depend only on the state at the end of the previous period and equations already solved.
With this in mind, also in the aggregate version we write the equations such that they can be solved in order, without interdependencies, to more easily adapt them to the AB version later.

Wages (W) are set exogenously and are the only input of the production, and so the only cost for the firms.

Prices (p) are so determined by a cost plus mark-up strategy like $$p = (1 + \mu) \frac{W}{Y}$$ where Y is the total output and $\mu$ is the mark-up rate.
By employing a Leontief production function the production ($\beta$) is the ratio between the single input and the output, and so prices can also be written as $$p= (1+\mu) \beta$$ which results in constant prices when productivity is constant.

Unemployment benefits (UB) are an exogenous and constant fraction $\phi$ of the wages.

Taxes (T) are paid with a flat rate on wages ($\tau_W$), distributed profits ($\tau_P$) and consumption ($\tau_C$).

Desired consumption in units of good for households is defined in a classic post-Keynesian way as $$C^T = \frac{a_Y ((1 - \tau_W) W + UB) + a_V {\mathbf{M}_\mathcal{H}}_{t-1}}{(1 + \tau_C) p_{t-1}}$$
where $a_Y$ and $a_V$ are the marginal propensity to consume out of disposable income (wages and unemployment benefits) and wealth (money, $\mathbf{M}_\mathcal{H}$).
Additionally, at the denominator appears $(1 + \tau_C) p_{t-1}$ which is the price of the good for the consumers (including VAT) and allows converting the desired expenditures for consumption into the desired units of goods.

Government's desired consumption is determined against target expected deficit $d$ $$G^T = \frac{d \cdot {GDP}_{t-1} + T_{t-1} - UB_{t-1}}{p_{t-1}}$$

Given the desired total consumption $Y^T = C^T + G^T$, firms aim to hire a number of workers equal to $N = Y^T/\beta$.
The number of households, or equivalently the number of available workers, is normalized in the $[0,1]$ interval.

Desired consumption are reduced proportionally to match the realized production $Y = \beta N$.

Finally, firms distributed any resulting money to households as profits (P) (or receive money from the households to cover the losses), netting their deposit to 0 at the end of each period.

The other equations simply mirror the sum of the rows and columns of the matrices of the model (\ref{tab:bs} and \ref{tab:fof}).

The exact implementation and the values used for the parameters and the initialization are reported in appendix.

\subsection{A comparison}
The model presented is close to the SIMEX model proposed by \textcite{godley2012}. This is the only step of our build up in which the two sequences of models are comparable, since their models focus on the financial and monetary side of the economy, while our models on the real side.

The two model are quite similar, sharing the same balance sheet, most of the flows and the equations, including a labour-only Leontief production function, a cost-plus-markup pricing and the standard post-Keynesian consumption function.
The flows' matrix has two more entries in our model: unemployment benefits and profits.

A first difference is the government spending which is endogenous in our model and exogenous in their one (and in the modified version presented in section \ref{sec:base_EC}).

A second one is about taxation: in Godley and Lavoie's model only income is taxed, while in our model also consumption and profits are getting taxed.

\subsection{Results}
\begin{figure}
    \centering
    \includegraphics[width=\textwidth]{00_sankey}
    \caption{Sankey diagram of flow at the steady state}
    \label{fig:00_san}
\end{figure}

Diagram \ref{fig:00_san} represents the flows of the model at the steady state.
The distribution of the flows (accounted by their nominal value rather than their real quantity) among the sector has the right sign and macroscopically realistic magnitudes\footnote{The only unusual value is the net transfer of money from the households to the government rather than the opposite: it is an artefact of the representation, the value of the associate flow is close to zero.}.

\begin{figure}
    \centering
    \begin{subfigure}[t]{0.49\textwidth}
        \centering
        \includegraphics[width=\textwidth]{00_M}
        \caption{Stock of money}
        \label{fig:00_M}
    \end{subfigure}%
    ~
    \begin{subfigure}[t]{0.49\textwidth}
        \centering
        \includegraphics[width=\textwidth]{00_N}
        \caption{Employment share}
        \label{fig:00_N}
    \end{subfigure}

    \begin{subfigure}[t]{0.49\textwidth}
        \centering
        \includegraphics[width=\textwidth]{00_Y}
        \caption{Output and desired output}
        \label{fig:00_Y}
    \end{subfigure}
    \caption{Dynamics of the base model}
\end{figure}

Figure \ref{fig:00_M} shows the stock of money in the simulation: as designed the firms save no money.

The model leads to full employment (figure \ref{fig:00_N}) with a potential output gap (figure \ref{fig:00_Y}).
In this model the government is allowed to influx liquidity each period (the target deficit is positive), so there are no financial constraints to the economy.
On the other side, given a fixed productivity and no population dynamics the maximum real output is fixed and is what prevent the economy to grow in real terms. Fixed wages and mark-up prevent a nominal growth.

\section{Exogenous government spending}
\label{sec:base_EC}
In a variation of the model (marked with \textit{EC} in the name of R script), we try to fix exogenously the government spending as considering $G^T$ as constant.

\begin{figure}
    \centering
    \includegraphics[width=\textwidth]{00_sankey_EC}
    \caption{Sankey diagram of flow at the steady state with exogenous government spending}
    \label{fig:00_san_EC}
\end{figure}

The Sankey diagram for this variant is shown in figure \ref{fig:00_san_EC}: the only macroscopic difference is the emergence of a non-zero value for the unemployment benefits, as a consequence of a non-zero unemployment rate (dependent of the level of government spending within a given range of stable values).

\begin{figure}
    \centering
    \includegraphics[width=.5\textwidth]{00_Y_EC}
    \caption{Output, desired output and (exogenous) government spending}
    \label{fig:00_Y_EC}
\end{figure}

The model dynamics is slower and no output gap appears (see figure \ref{fig:00_Y_EC}): the exogenous government spending limits the influx of money in the system; as a consequence the cash flow of the households is limited and their consumption as well.
In such a scenario the goods market is short on the demand side rather than on the supply side and, as a consequence, not all the production capacity of the system is required, leading to unemployment.
The choice to make good perishable and not keep them into inventories prevent the model to develop explosive cyclical dynamics in which the good are stocked into the inventories which are later sold interrupting the production and reducing the employment, leading to a further reduction of the demand.

\section{The agent-based extension}
To adapt the aggregate model to an Agent-Based one there are few choices to be made about the distribution and the discretization of some variables.

First, the households and firms sectors are now a set of agents rather than being represented as a single variable.
Second, the non-zero values at the beginning of the model (in our case the stock of money held by the households) have to be divided among the agent.
Third, each interaction among non-aggregate sectors now requires a matching procedure.

As stated at the beginning of section \ref{sec:behave}, in an ABM actions have to happen in a specific order, to ensure that there are not decision that agents have to take relies on of future actions or choices.
The sequence of the events of the model is:
\begin{enumerate}[label=\textbf{\Alph*}]
    \item Households determinate the desired consumption (expressed as units of good) based on the wealth and the disposable income of the previous period
    \item Government determinate the desired consumption (expressed as units of good) based on the target deficit computed on the accounting of the previous period
    \item Firms fire and hire workers in order to match the labour input required to match the desired consumption
    \item The realized output is determined
    \item Wages and unemployment benefits are paid
    \item Prices are set
    \item The goods are sold to households and the government
    \item Profits are distributed
    \item Taxes are paid
    \item Accounting variables are updated
\end{enumerate}

The model is initialized by attributing the same amount of money to each household: in this way is easier to observe if and how inequalities emerge.

All the matching procedures present a strong path-dependency to reduce the computational cost: each period the resulting matching of the previous one is refined instead to be reconstructed from scratch (as if all the working contracts were permanent and the good market show a kind of monopolistic competition that sticks together producers and consumers).

The job market matching is random: if the aggregate desired consumption by both the households and the government is higher than the output related to the current level of employment a random firm employs a random unemployed household; otherwise if the current level of employment overshoots the target demand a random firm fires a random household.

To distribute the consumption between the households and the public sector, the two sectors get the same fraction of the desired consumption such that all the produced goods are bought (or the entire desired consumption is matched and the over-production is lost). The unsatisfied consumption of the households is distributed randomly among them (or, in fact, the matching for the consumption goods is random).

Profits are distributed from all firms to all households proportionally to the stock of money held by each household at the end of the previous period. This resembles a homogenous return on wealth, as in modern ETFs.

\subsection{Results}
\begin{figure}
    \centering
    \begin{subfigure}[t]{0.49\textwidth}
        \centering
        \includegraphics[width=\textwidth]{00AB_M}
        \caption{Stock of money and net value of the sectors}
        \label{fig:00AB_M}
    \end{subfigure}%
    ~
    \begin{subfigure}[t]{0.49\textwidth}
        \centering
        \includegraphics[width=\textwidth]{00AB_macro}
        \caption{Consumption, output and employment share (dotted is desired)}
        \label{fig:00AB_macro}
    \end{subfigure}

    \begin{subfigure}[t]{0.49\textwidth}
        \centering
        \includegraphics[width=\textwidth]{00AB_moments}
        \caption{Moments of wealth distribution among households}
        \label{fig:00AB_moments}
    \end{subfigure}%
    ~
    \begin{subfigure}[t]{0.49\textwidth}
        \centering
        \includegraphics[width=\textwidth]{00AB_hist}
        \caption{Evolution of wealth distribution among households}
        \label{fig:00AB_hist}
    \end{subfigure}
    \caption{Dynamics of AB base model}
\end{figure}

The qualitative dynamics of the stock and flows of the model remains largely the same: the economy is still at full employment with unsatisfied demand (figure \ref{fig:00AB_macro}).
The AB extension introduces some volatility in the macro variables as consequence of the heterogeneity among agents: figure \ref{fig:00AB_moments} shows how inequality among agents varies during the simulation introducing some form of cycles in the dynamics.

Finally, even such a simple model is able to produce a fat tail distribution of wealth (see figure \ref{fig:00AB_hist}): this is not a guarantee result since many previous versions in which all the households demand was fulfilled before the public one, were not able to get such a dynamics for the effects of the central limit theorem.
Introducing path dependency and limited consumption allow for some agents to under consume (and so having a higher saving rate) and to accumulate an initial stock of money. The proportional distribution of profits amplifies this dynamics, but it is not necessary: even by dividing the profits equally among agents, the same initial inequality appears.

\subsection{Matrix-based vs OOP comparison}
In this version of the model (and in the next one) the qualitative behaviour of the two implementations (matrix-based in R and OOP-based in Python) is almost identical.
The OOP version is significantly slower, because the agents are simulated one-by-one, rather than all together as a single row in the matrix.
This allows for more complex behavioural equations and interactions at the cost, as said, of increased computational time and resources.

Another difference to be noted is that the matrix-based approach allows for updating the stocks at the end of the period, while the OOP version updates them continuously.
This highlights the double and quadruple entries' principle: all the stocks and the flows must be updated in both the agents involved, which results in two lines of code each updating the information for one of the two.
Moreover, each flow requires the update of two stocks, one for each agent involved.

The graphs above are taken from the R version.